%%%%%%%%%%  scan idx 225  |  folio 212  %%%%%%%%%%
\emph{Preuve.} Il est clair que $u_X$ est continue. D'autre part (Bourbaki, Intégration, chap. III \S~2, lemmes 1 et 2), le sous-espace $C(X,\mathbb C)\otimes_{\mathbb C}E$ de $C(X,E)$ (formé des applications continues de $X$ dans $E$ qui prennent leurs valeurs dans un sous-espace vectoriel de dimension finie) est dense~; de même $C(X,\mathbb C)\otimes_{\mathbb C}F$ est dense dans $C(X,F)$. Supposons $u$ continue et d'image dense. Alors
\[
  u_X(C(X,\mathbb C)\otimes_{\mathbb C}E)=C(X,\mathbb C)\otimes_{\mathbb C}u(E)
\]
est dense dans $C(X,\mathbb C)\otimes_{\mathbb C}F$, d'où l'on déduit, en prenant les adhérences, que $u_X(C(X,E))$ est dense dans $C(X,F)$, cqfd.

\textbf{3.1.3.} Soient $V$ une variété compacte, $E$ et $F$ des fibrés vectoriels complexes\footnote{Pour fixer les idées (résultats analogues dans le cas réel).} de rang fini sur $V$, et $n\in\N$. L'espace $\Diff^n_\mathbb C(E,F)$, s'identifiant par (1.2.3.1) à
\[
  \Homop(P^n(E),F)\simeq C^\infty(V,\underline{\Homop}(P^n(E),F)),
\]
est muni naturellement d'une structure d'espace de Fréchet. Soit $u\in\Diff^n_\mathbb C(E,F)$~; on vérifie facilement que l'application linéaire de $C^\infty(V,E)$ dans $C^\infty(V,F)$ définie par $u$ est continue. Si $\Lop(C^\infty(V,E),C^\infty(V,F))$ désigne l'espace des applications linéaires continues de $C^\infty(V,E)$ dans $C^\infty(V,F)$, on a donc une application linéaire
\[
  \varphi:\Diff^n_\mathbb C(E,F)
  \longrightarrow \Lop(C^\infty(V,E),C^\infty(V,F)).
\]
D'autre part, pour $r\in\Z$, $u$ agit par extension des scalaires ([16]~IX Th.~12) sur les sections $C^r_2$ de $E$ et définit une application linéaire
\[
  \varphi_r:\Diff^n_\mathbb C(E,F)
  \longrightarrow \Lop(C^r_2(V,E),C^{r-n}_2(V,F)).
\]
On vérifie aisément que~:
\begin{enumerate}[label=(\roman*)]
\item les $\varphi_r$ sont \emph{continues} (l'espace d'arrivée est un espace de Banach 
